\documentclass[12pt,a4paper]{article}

% Drake_preamble.tex - LaTeX Preamble Template
% 作者: 謝慕揚, MD, PhD, FESC
% 
% 使用說明:
% 1. 表格請使用 longtable，不要有垂直分隔線 |
% 2. 表格內換行使用 \newline，不要使用 \\
% 3. 藥物名稱、檢驗、檢查請維持英文
% 4. Reference 格式請使用 \href 加上驗證過的 DOI

% Including essential packages for Chinese typesetting
\usepackage{xeCJK}
\usepackage{fontspec}
% Configuring Chinese font with Noto Serif CJK TC, fallback to SimSun
\IfFontExistsTF{Noto Serif CJK TC}{
  \setCJKmainfont{Noto Serif CJK TC}
  \setCJKsansfont{Noto Serif CJK TC}
  \setCJKmonofont{Noto Serif CJK TC}
}{\setCJKmainfont{SimSun}
  \setCJKsansfont{SimSun}
  \setCJKmonofont{SimSun}
}

% Including other necessary packages
\usepackage{geometry}
\usepackage{titlesec}
\usepackage{enumitem}
\usepackage{xcolor}
\usepackage{colortbl}
\usepackage{tcolorbox}
\usepackage{booktabs}
\usepackage{longtable}
\usepackage{array}
\usepackage{graphicx}
\usepackage{tikz}
\usepackage{fancyhdr}
\usepackage{multicol}
\usepackage{datetime2}
\usepackage{hyperref}
\usepackage{mdframed}
\usepackage{amsmath}
\usepackage{amssymb}
\usepackage{multirow}

\hypersetup{
    colorlinks=true,
    linkcolor=emergency_blue,
    citecolor=emergency_blue,
    urlcolor=emergency_blue,
    pdfborder={0 0 0}
}

% Defining page geometry
\geometry{
  top=2.5cm,
  bottom=2.5cm,
  left=2cm,
  right=2cm,
  headheight=25pt
}

% Adjusting topmargin to compensate for increased headheight
\addtolength{\topmargin}{-10pt}

% Defining colors
\definecolor{emergency_red}{RGB}{186,24,27}
\definecolor{emergency_blue}{RGB}{0,114,188}
\definecolor{emergency_gray}{RGB}{120,120,120}
\definecolor{highlight_yellow}{RGB}{255,250,205}

% Setting title formats
\titleformat{\section}[block]{\Large\bfseries\color{emergency_red}}{\thesection}{10pt}{}
\titleformat{\subsection}[block]{\large\bfseries\color{emergency_blue}}{\thesubsection}{8pt}{}
\titleformat{\subsubsection}[block]{\normalsize\bfseries\color{emergency_gray}}{\thesubsubsection}{6pt}{}

% Defining custom environment for important notes
\newmdenv[
    linecolor=emergency_red,
    backgroundcolor=highlight_yellow,
    linewidth=2pt,
    topline=true,
    bottomline=true,
    leftline=true,
    rightline=true,
    innertopmargin=10pt,
    innerbottommargin=10pt,
    innerleftmargin=10pt,
    innerrightmargin=10pt
]{importantbox}

% Defining note command with tcolorbox
\newcommand{\note}[1]{
    \par
    \noindent
    \begin{tcolorbox}[
        colback=emergency_blue!5!white,
        colframe=emergency_blue,
        title=\textbf{重點提醒},
        fonttitle=\bfseries,
        boxrule=0.5pt,
        arc=3pt,
        left=6pt,
        right=6pt,
        top=6pt,
        bottom=6pt
    ]
    #1
    \end{tcolorbox}
}

% Key findings box
\newcommand{\keyfinding}[1]{
    \par
    \noindent
    \begin{tcolorbox}[
        colback=yellow!10!white,
        colframe=orange!75!black,
        title=\textbf{核心發現摘要},
        fonttitle=\bfseries,
        boxrule=0.5pt,
        arc=3pt
    ]
    #1
    \end{tcolorbox}
}

% Defining custom date format for ISO 8601
\DTMnewdatestyle{iso}{
  \renewcommand{\DTMdisplaydate}[4]{\number##1-\number##2-\number##3}
  \renewcommand{\DTMDisplaydate}[4]{\number##1-\number##2-\number##3}
}
\DTMsetdatestyle{iso}
\newcommand{\mydate}{\DTMtoday}


% Custom header for this document
\pagestyle{fancy}
\fancyhf{}
\fancyhead[L]{\textcolor{emergency_red}{NEJM Clinical Practice 教學講義}}
\fancyhead[R]{\textcolor{emergency_blue}{Peripheral Artery Disease}}
\fancyfoot[C]{\textcolor{emergency_gray}{\thepage}}
\renewcommand{\headrulewidth}{1pt}
\renewcommand{\headrule}{\hbox to\headwidth{\color{emergency_red}\leaders\hrule height \headrulewidth\hfill}}

\begin{document}

%============================================================
% Title Page
%============================================================
\begin{center}
{\LARGE\bfseries\textcolor{emergency_red}{New England Journal of Medicine}}\\[0.3cm]
{\Large\textcolor{emergency_blue}{Clinical Practice}}\\[0.5cm]
{\Large\textbf{Peripheral Artery Disease in the Legs}}\\[0.3cm]
{\large 下肢周邊動脈疾病}\\[0.5cm]
{\normalsize \href{https://doi.org/10.1056/NEJMcp2501200}{N Engl J Med 2026;394:486-96. DOI: 10.1056/NEJMcp2501200}}\\[0.3cm]
{\normalsize 作者：Mary M. McDermott, MD (Northwestern University)}\\[0.3cm]
{\normalsize 整理：謝慕揚 MD, PhD, FESC \quad 日期：\mydate}
\end{center}

\vspace{0.5cm}

%============================================================
\section{案例導讀 (Clinical Vignette)}
%============================================================

一位 62 歲女性，有 type 2 diabetes、hypertension 及吸菸史，主訴行走時出現下肢症狀。她描述行走時臀部及小腿疼痛，休息 10 分鐘內可緩解。因症狀限制，步行超過一至兩個街區即感困難。

\textbf{理學檢查：}
\begin{itemize}[leftmargin=2em]
    \item 血壓 150/70 mmHg，心率 82 bpm
    \item 無 femoral bruits
    \item Calf muscles atrophy
    \item Dorsalis pedis 及 posterior tibial pulses 雙側均為 grade 1+ (正常為 3+)
    \item \textbf{Ankle-brachial index (ABI)：右側 0.56，左側 0.64}
\end{itemize}

\note{此案例為典型的 peripheral artery disease (PAD) 合併 claudication，且 ABI 顯示中度至重度疾病。本文將探討如何診斷、治療行走障礙，以及預防心血管事件。}

%============================================================
\section{流行病學 (Epidemiology)}
%============================================================

\subsection{盛行率}

\begin{itemize}[leftmargin=2em]
    \item 全球約影響 \textbf{2.36 億人}
    \item 美國約 \textbf{1,250 萬人} ($\geq$40 歲)
    \item 65 歲以上：10-15\%
    \item 80 歲以上：15-20\%
    \item 美國 Black 族群盛行率約為 non-Hispanic White 的 \textbf{2 倍}
\end{itemize}

\subsection{危險因子 (Risk Factors)}

\begin{longtable}{p{4cm}p{10cm}}
\toprule
\textbf{危險因子} & \textbf{說明} \\
\midrule
\endfirsthead
\toprule
\textbf{危險因子} & \textbf{說明} \\
\midrule
\endhead
\bottomrule
\endfoot
年齡 & 年齡越大風險越高 \\
\midrule
吸菸 (Cigarette smoking) & 最重要的可改變危險因子 \\
\midrule
糖尿病 (Diabetes mellitus) & PAD 合併 DM 者有更嚴重的行走障礙、更高的死亡率、截肢率，且 below-the-knee atherosclerosis 更常見 \\
\midrule
血脂異常 (Dyslipidemia) & 包括 elevated lipoprotein(a) \\
\midrule
高血壓 (Hypertension) & 加速動脈粥狀硬化 \\
\midrule
久坐生活型態 & 缺乏運動 \\
\bottomrule
\end{longtable}

\subsection{不良預後 (Adverse Outcomes)}

\begin{importantbox}
\textbf{心血管事件風險}
\begin{itemize}[leftmargin=2em]
    \item PAD 患者心血管事件及死亡率為無 PAD 者的 \textbf{2-3 倍}
    \item 10 年心血管死亡率：
    \begin{itemize}
        \item 男性：PAD 18.7\% vs. 無 PAD 4.4\%
        \item 女性：PAD 12.6\% vs. 無 PAD 4.1\%
    \end{itemize}
    \item 症狀性 PAD (ABI <0.80 或曾接受 revascularization) 患者，30 個月追蹤期間有 \textbf{10.7\%} 發生 MI、ischemic stroke 或心血管死亡
\end{itemize}
\end{importantbox}

\textbf{Major adverse limb events (MALE):}
\begin{itemize}[leftmargin=2em]
    \item 定義：severe leg ischemia 導致 revascularization 或 amputation
    \item 發生率：1.8-2.7 年追蹤期間約 12.9-15.2\%
    \item 接受 leg revascularization 後更常見
\end{itemize}

%============================================================
\section{臨床表現與診斷 (Clinical Presentation \& Diagnosis)}
%============================================================

\subsection{症狀表現}

\keyfinding{
\textbf{重要發現：}高達 68-90\% 的 PAD 患者\textbf{沒有典型 claudication 症狀}或完全無症狀！
\begin{itemize}[leftmargin=2em]
    \item 典型 claudication (classic)：僅約 10-32\%
    \item 非典型症狀 (atypical)：如髖部疼痛、下背痛
    \item 無症狀 (asymptomatic)：約 58.5\%（德國研究）
\end{itemize}
許多患者會減少行走活動或速度以避免缺血症狀，這可能導致\textbf{診斷不足}。
}

\textbf{Classic intermittent claudication 特徵：}
\begin{itemize}[leftmargin=2em]
    \item 運動時小腿疼痛
    \item 休息時不會發生
    \item 休息 10 分鐘內可緩解
    \item 行走距離越長症狀越嚴重
\end{itemize}

\subsection{理學檢查}

脈搏觸診對診斷 PAD 的準確度有限：

\begin{longtable}{p{5cm}p{3cm}p{3cm}}
\toprule
\textbf{檢查項目} & \textbf{Sensitivity} & \textbf{Specificity} \\
\midrule
\endfirsthead
Absent dorsalis pedis pulse & 50\% & 73\% \\
\midrule
Absent posterior tibial pulse & 71\% & 91\% \\
\bottomrule
\end{longtable}

\subsection{Ankle-Brachial Index (ABI)}

\textbf{定義：}踝部 systolic pressure (posterior tibial 或 dorsalis pedis) 與 brachial artery systolic pressure 的比值，使用 Doppler ultrasonography 測量。

\begin{longtable}{p{4cm}p{10cm}}
\toprule
\textbf{ABI 數值} & \textbf{意義} \\
\midrule
\endfirsthead
>1.00 - 1.40 & 正常 \\
\midrule
<0.90 & 診斷 PAD (>50\% angiographic stenosis)\newline Sensitivity 69-79\%, Specificity 83-99\% \\
\midrule
0.70 - <0.90 & 輕度 PAD (mild) \\
\midrule
0.50 - <0.70 & 中度 PAD (moderate) \\
\midrule
<0.50 & 重度 PAD (severe) \\
\midrule
$\geq$1.40 & 疑似 medial calcification (不可壓縮動脈)\newline 常見於 diabetes 或 CKD 患者 \\
\bottomrule
\end{longtable}

\note{
\textbf{Toe-brachial index (TBI)：}當 ABI $\geq$1.40 或症狀明顯但 ABI 正常時，可使用 TBI。因趾動脈較不易發生 medial calcification。\textbf{TBI $\leq$0.70} 表示 PAD。
}

\subsection{行走能力與疾病嚴重度}

\begin{longtable}{p{4cm}p{6cm}p{4cm}}
\toprule
\textbf{ABI} & \textbf{疾病嚴重度} & \textbf{無法完成 6 分鐘步行測試\newline 需中途休息比例} \\
\midrule
\endfirsthead
1.10 - 1.40 & 正常 & 3.3\% \\
\midrule
0.70 - <0.90 & 輕度 (mild) & 18.1\% \\
\midrule
0.50 - <0.70 & 中度 (moderate) & 31.6\% \\
\midrule
<0.50 & 重度 (severe) & 52.0\% \\
\bottomrule
\end{longtable}

PAD 患者的行走能力會隨時間下降：
\begin{itemize}[leftmargin=2em]
    \item 6 個月：6-minute walk distance 平均減少 10.2 m
    \item 18 個月：6-minute walk distance 平均減少 34 m
\end{itemize}

%============================================================
\section{治療：改善行走障礙}
%============================================================

\subsection{治療效果比較圖}

\begin{center}
\begin{tabular}{p{6cm}p{4cm}p{4cm}}
\toprule
\textbf{治療方式} & \textbf{Maximal treadmill\newline walking distance 增加} & \textbf{6-minute walk\newline distance 增加} \\
\midrule
Cilostazol & +40 m & -- \\
\midrule
Semaglutide & +39.9 m & -- \\
\midrule
Liraglutide & -- & +25 m \\
\midrule
Home-based walking exercise & +53.7 m & +55.6 m \\
\midrule
Femoropopliteal endovascular & +69 m & -- \\
\midrule
Supervised treadmill exercise & +186 m & +31.8 m \\
\midrule
Aortoiliac endovascular & +202 m & -- \\
\bottomrule
\end{tabular}
\end{center}

\subsection{藥物治療}

\subsubsection{Cilostazol (Phosphodiesterase 3 inhibitor)}

\begin{itemize}[leftmargin=2em]
    \item \textbf{適應症：}PAD 導致的 claudication 及行走障礙
    \item \textbf{劑量：}起始 50 mg BID，可調整至 100 mg BID
    \item \textbf{效果：}Maximal treadmill walking distance 增加約 40 m（vs. placebo）
    \item \textbf{副作用：}Headache, palpitations, diarrhea
    \item \textbf{禁忌症：}\textcolor{emergency_red}{\textbf{任何程度的 heart failure}}（其他 PDE3 inhibitors 會增加 HF 患者死亡率）
    \item \textbf{建議：}若 12 週後行走能力無改善，應停藥
\end{itemize}

\subsubsection{GLP-1 Receptor Agonists}

\textbf{Semaglutide (STRIDE trial):}
\begin{itemize}[leftmargin=2em]
    \item 792 位 PAD + T2DM 患者
    \item 治療 52 週
    \item Maximal treadmill walking distance ratio: 1.13 (95\% CI 1.06-1.21, P<0.001)
    \item 絕對增加約 39.9 m
\end{itemize}

\textbf{Liraglutide:}
\begin{itemize}[leftmargin=2em]
    \item 55 位 PAD + T2DM 患者（open-label）
    \item 6 個月後 6-minute walk distance 增加 25 m (P<0.001)
\end{itemize}

\subsection{運動治療 (Exercise Therapy)}

\begin{importantbox}
\textbf{AHA/ACC Guidelines 建議：}Supervised walking exercise 及 structured home-based walking exercise 均為\textbf{第一線治療}。
\end{importantbox}

\subsubsection{Supervised Walking Exercise}

\begin{longtable}{p{4cm}p{10cm}}
\toprule
\textbf{項目} & \textbf{說明} \\
\midrule
\endfirsthead
定義 & 在醫療機構由 nurse 或 exercise physiologist 監督下進行 treadmill walking \\
\midrule
頻率 & 每週 3 次 \\
\midrule
目標 & 每次運動誘發最大缺血性腿部症狀，逐步增加至每次 50 分鐘 \\
\midrule
效果 & Maximal treadmill walking distance +180 m\newline 6-minute walk distance +31.8 m \\
\midrule
保險給付 & Medicare 給付 36 次/12 週（症狀性 PAD）\newline 可涵蓋 walking、cycling、resistance training \\
\midrule
注意事項 & 約 6 週後開始有明顯效果\newline 停止運動後效果消失 \\
\bottomrule
\end{longtable}

\subsubsection{Structured Home-based Walking Exercise}

\begin{longtable}{p{4cm}p{10cm}}
\toprule
\textbf{項目} & \textbf{說明} \\
\midrule
\endfirsthead
定義 & 在家中或居家附近進行非監督式步行運動，但有 coach 透過電話或定期面談追蹤 \\
\midrule
頻率 & 每週 3-5 次，每次 30-50 分鐘 \\
\midrule
強度 & \textcolor{emergency_red}{\textbf{必須誘發缺血性腿部症狀 (high intensity)}} \newline 低強度運動無顯著效果 \\
\midrule
效果 & Maximal treadmill walking distance +53.7 m\newline 6-minute walk distance +55.6 m \\
\midrule
行為介入 & 設定運動目標\newline 管理運動時缺血性疼痛的策略\newline 監測進度\newline 建立 self-efficacy \\
\midrule
保險給付 & Medicare 不給付 \\
\midrule
優點 & 方便、可近性高 \\
\bottomrule
\end{longtable}

\note{
\textbf{LITE trial 重要發現：}High-intensity home-based exercise (誘發缺血症狀) 使 6-minute walk distance 增加 49.6 m (P<0.001)，而 low-intensity exercise (不誘發症狀) 僅增加 8.7 m，與對照組無顯著差異 (P=0.44)。
}

\subsection{血管重建術 (Revascularization)}

\begin{importantbox}
\textbf{適應症：}僅限於接受運動治療後仍有\textbf{持續性、影響生活品質的缺血性腿部症狀}之患者。
\end{importantbox}

\subsubsection{手術方式}

\begin{itemize}[leftmargin=2em]
    \item \textbf{Surgical revascularization:} Endarterectomy 或 bypass surgery
    \item \textbf{Endovascular revascularization:}
    \begin{itemize}
        \item Uncoated 或 drug-coated balloon angioplasty
        \item Bare-metal 或 drug-eluting stents
        \item Atherectomy 或 lithotripsy（處理嚴重鈣化病灶）
    \end{itemize}
\end{itemize}

\subsubsection{效果比較}

\begin{itemize}[leftmargin=2em]
    \item \textbf{Aortoiliac endovascular:} 效果最佳且最持久，與 supervised exercise 相當
    \item \textbf{Femoropopliteal endovascular:} 效果較差
    \item \textbf{Infrapopliteal arteries:} 對 claudication 患者效果不明確
\end{itemize}

\textbf{CLEVER trial (Aortoiliac PAD):}
\begin{itemize}[leftmargin=2em]
    \item Maximal treadmill walking time：supervised exercise = endovascular revascularization
    \item Patient-reported outcomes (Walking Impairment Questionnaire)：\newline Endovascular > Supervised exercise (score change 43.8 vs. 25.1, P=0.03)
\end{itemize}

\subsubsection{Endovascular vs. Open Bypass Surgery}

\begin{longtable}{p{5cm}p{4cm}p{4cm}}
\toprule
\textbf{項目} & \textbf{Endovascular} & \textbf{Open Bypass} \\
\midrule
\endfirsthead
Repeat revascularization (1 年) & 12\% & 5\% \\
\midrule
Repeat revascularization (3 年) & 19\% & 8\% \\
\midrule
30 天 procedure-related complications & 8.8\% & 30.7\% \\
\bottomrule
\end{longtable}

%============================================================
\section{心血管事件預防}
%============================================================

\subsection{Cholesterol-lowering Therapy}

\begin{itemize}[leftmargin=2em]
    \item \textbf{建議：}High-intensity statin (atorvastatin 或 rosuvastatin)，使用可耐受的最高劑量
    \item \textbf{目標：}LDL-C 降低至少 50\%（ESC 建議 <55 mg/dL）
    \item \textbf{證據：}Heart Protection Study - Simvastatin 40 mg vs. placebo
    \begin{itemize}
        \item 全因死亡：12.9\% vs. 14.7\% (P<0.001)
        \item 冠心病死亡：5.7\% vs. 6.9\% (P<0.001)
    \end{itemize}
    \item \textbf{加強治療：}高風險患者可加用 PCSK9 inhibitor 或 ezetimibe，可進一步減少 CV events 及 MALE
\end{itemize}

\subsection{Antiplatelet and Antithrombotic Therapy}

\begin{longtable}{p{4cm}p{10cm}}
\toprule
\textbf{治療選項} & \textbf{說明} \\
\midrule
\endfirsthead
Single antiplatelet & Aspirin 81-325 mg/day 或 Clopidogrel 75 mg/day \\
\midrule
Dual pathway inhibition & \textbf{Rivaroxaban 2.5 mg BID + Aspirin 100 mg/day} \\
\bottomrule
\end{longtable}

\textbf{COMPASS trial (PAD or carotid artery disease, n=7,470):}

\begin{center}
\begin{tabular}{p{5cm}p{3cm}p{3cm}p{2cm}}
\toprule
\textbf{Outcome} & \textbf{Rivaroxaban + ASA} & \textbf{ASA alone} & \textbf{P value} \\
\midrule
Cardiovascular events & 5\% & 7\% & 0.005 \\
Major adverse limb events & 1\% & 2\% & 0.004 \\
Major bleeding & 3\% & 2\% & -- \\
\bottomrule
\end{tabular}
\end{center}

\textbf{VOYAGER PAD trial (post-revascularization, n=6,564):}
\begin{itemize}[leftmargin=2em]
    \item Rivaroxaban 2.5 mg BID + ASA vs. ASA alone
    \item 3 年 composite outcome (CV death, MI, stroke, ALI, major amputation)：17.3\% vs. 19.9\% (P=0.009)
    \item Major bleeding 較高
\end{itemize}

\subsection{Blood Pressure Treatment}

\begin{itemize}[leftmargin=2em]
    \item \textbf{目標：}<130/80 mmHg
    \item \textbf{證據：}SPRINT trial - 高風險患者降至 <120 mmHg vs. <140 mmHg
    \begin{itemize}
        \item CV events：1.77\%/year vs. 2.40\%/year
    \end{itemize}
\end{itemize}

\subsection{GLP-1 Receptor Agonists (CV Prevention)}

\textbf{SELECT trial (CV disease + BMI $\geq$27, no DM, n=17,604):}
\begin{itemize}[leftmargin=2em]
    \item 8.7\% 有症狀性 PAD
    \item Semaglutide vs. placebo
    \item CV events：6.5\% vs. 8.0\% (P<0.001)，追蹤 39.8 個月
\end{itemize}

\textbf{SOUL trial (T2DM + ASCVD or CKD, n=9,650):}
\begin{itemize}[leftmargin=2em]
    \item 15.7\% 有 PAD
    \item Oral semaglutide 14 mg vs. placebo
    \item CV events：12.0\% vs. 13.8\% (P=0.006)，追蹤 47.5 個月
\end{itemize}

\subsection{SGLT2 Inhibitors (糖尿病患者)}

\begin{longtable}{p{3cm}p{5cm}p{6cm}}
\toprule
\textbf{Trial} & \textbf{藥物/族群} & \textbf{結果} \\
\midrule
\endfirsthead
EMPA-REG & Empagliflozin\newline T2DM 高 CV 風險\newline 21\% 有 PAD & CV events: 10.5\% vs. 12.1\% (P=0.04)\newline 無截肢風險增加 \\
\midrule
CANVAS & Canagliflozin\newline T2DM 高 CV 風險\newline 21\% 有 PVD & CV events 減少 (P=0.02)\newline \textcolor{emergency_red}{截肢風險增加 (6.3 vs. 3.4/1000 pt-yr)} \\
\midrule
CREDENCE & Canagliflozin\newline T2DM + CKD\newline 23.8\% 有 PVD & 截肢風險無顯著增加 (12.3 vs. 11.2/1000 pt-yr) \\
\bottomrule
\end{longtable}

\note{
\textbf{臨床建議：}對 PAD + T2DM 患者，\textbf{empagliflozin 或 dapagliflozin} 可能是較安全的選擇，因未觀察到截肢風險增加。
}

\subsection{Smoking Cessation}

\begin{itemize}[leftmargin=2em]
    \item 使用行為介入 + 藥物 (varenicline, bupropion, 或 nicotine replacement)
    \item \textbf{證據：}PAD 患者接受 leg angiography 後 1 年
    \begin{itemize}
        \item 戒菸者 vs. 持續吸菸者
        \item 全因死亡：14\% vs. 31\%
        \item Amputation-free survival：81\% vs. 60\%
    \end{itemize}
\end{itemize}

%============================================================
\section{治療摘要與指引建議}
%============================================================

\subsection{AHA/ACC 2024 Guidelines 建議}

\begin{longtable}{p{5cm}p{9cm}}
\toprule
\textbf{目標} & \textbf{建議} \\
\midrule
\endfirsthead
\toprule
\textbf{目標} & \textbf{建議} \\
\midrule
\endhead
\bottomrule
\endfoot
預防心血管事件 & Intensive cholesterol-lowering (statin)\newline BP <130/80 mmHg\newline Single antiplatelet 或 low-dose rivaroxaban + aspirin \\
\midrule
預防 MALE & Intensive cholesterol-lowering\newline Low-dose rivaroxaban + aspirin \\
\midrule
改善行走障礙 & Supervised walking exercise (first-line)\newline Structured home-based walking exercise (first-line)\newline Cilostazol (心衰竭禁用)\newline Revascularization (運動治療無效時) \\
\bottomrule
\end{longtable}

\subsection{ESC 2024 Guidelines 差異}

\begin{itemize}[leftmargin=2em]
    \item LDL-C 目標：<55 mg/dL（較 AHA/ACC 更嚴格）
    \item Home-based exercise：視為「reasonable」但非第一線治療
\end{itemize}

%============================================================
\section{Key Points 重點整理}
%============================================================

\begin{importantbox}
\begin{enumerate}[leftmargin=2em]
    \item 高達 90\% 的 PAD 患者\textbf{沒有典型 claudication}，約 60\% 無症狀

    \item \textbf{ABI <0.90} 診斷 PAD 的準確度為 72-89\%

    \item \textbf{10 年心血管死亡率}：PAD 患者顯著高於無 PAD 者
    \begin{itemize}
        \item 男性：18.7\% vs. 4.4\%
        \item 女性：12.6\% vs. 4.1\%
    \end{itemize}

    \item \textbf{6 分鐘步行測試}需中途休息比例：
    \begin{itemize}
        \item 無 PAD：3.3\%
        \item 輕度 PAD：18.1\%
        \item 重度 PAD：52.0\%
    \end{itemize}

    \item \textbf{心血管預防}：Statin + 抗血小板/抗凝血 + 血壓控制 + GLP-1 RA + SGLT2i (DM)

    \item \textbf{行走障礙治療}：
    \begin{itemize}
        \item 運動治療為第一線 (supervised 或 home-based)
        \item Cilostazol、semaglutide 效果有限
        \item Revascularization 僅用於運動治療無效者
    \end{itemize}
\end{enumerate}
\end{importantbox}

%============================================================
\section{案例處理建議}
%============================================================

對於本文開頭的 62 歲女性案例，作者建議：

\textbf{改善行走障礙：}
\begin{enumerate}[leftmargin=2em]
    \item 若能配合每週三次就診且有保險給付 → \textbf{Supervised exercise 12 週}
    \item 若無法配合 supervised exercise → \textbf{Structured home-based walking exercise}
    \item \textbf{Cilostazol 或 semaglutide} 可作為輔助治療
    \item 若以上治療仍無法改善症狀 → 考慮 \textbf{revascularization}
\end{enumerate}

\textbf{預防心血管事件：}
\begin{itemize}[leftmargin=2em]
    \item \textbf{Atorvastatin 或 rosuvastatin}：可耐受的最高劑量
    \item \textbf{血壓控制}：<130/80 mmHg
    \item \textbf{抗血栓治療}：Rivaroxaban 2.5 mg BID + low-dose aspirin，或 clopidogrel alone
    \item 考慮 \textbf{SGLT2 inhibitor + semaglutide}（減少心血管事件）
\end{itemize}

%============================================================
\section{未解問題 (Areas of Uncertainty)}
%============================================================

\begin{enumerate}[leftmargin=2em]
    \item 各種治療（運動、藥物、revascularization）的\textbf{最佳組合與時機}尚待確立
    \item \textbf{抗發炎藥物}（如 canakinumab）對行走能力的效果需進一步研究
    \item \textbf{促進骨骼肌生長藥物}（如 bimagrumab）的效果待研究
    \item \textbf{Colchicine} 等抗發炎藥物預防 PAD 患者心血管事件的效果尚不明確
\end{enumerate}

%============================================================
\section{參考文獻}
%============================================================

\begin{enumerate}
    \item McDermott MM. Peripheral Artery Disease in the Legs. \href{https://doi.org/10.1056/NEJMcp2501200}{\textit{N Engl J Med}. 2026;394:486-96.}

    \item Gornik HL, Aronow HD, Goodney PP, et al. 2024 ACC/AHA/AACVPR/APMA/ABC/SCAI/SVM/SVN/SVS/SIR/VESS Guideline for the Management of Lower Extremity Peripheral Artery Disease. \href{https://doi.org/10.1161/CIR.0000000000001251}{\textit{Circulation}. 2024;149:e1313-e1410.}

    \item Mazzolai L, Teixido-Tura G, Lanzi S, et al. 2024 ESC Guidelines for the Management of Peripheral Arterial and Aortic Diseases. \href{https://doi.org/10.1093/eurheartj/ehae179}{\textit{Eur Heart J}. 2024;45:3538-700.}

    \item Bonaca MP, Catarig A-M, Houlind K, et al. Semaglutide and Walking Capacity in People with Symptomatic Peripheral Artery Disease and Type 2 Diabetes (STRIDE). \href{https://doi.org/10.1016/S0140-6736(25)00565-9}{\textit{Lancet}. 2025;405:1580-93.}

    \item Anand SS, Bosch J, Eikelboom JW, et al. Rivaroxaban with or without Aspirin in Patients with Stable Peripheral or Carotid Artery Disease (COMPASS). \href{https://doi.org/10.1016/S0140-6736(17)32409-1}{\textit{Lancet}. 2018;391:219-29.}

    \item Bonaca MP, Bauersachs RM, Anand SS, et al. Rivaroxaban in Peripheral Artery Disease after Revascularization (VOYAGER PAD). \href{https://doi.org/10.1056/NEJMoa1910459}{\textit{N Engl J Med}. 2020;382:1994-2004.}

    \item McDermott MM, Spring B, Tian L, et al. Effect of Low-Intensity vs High-Intensity Home-Based Walking Exercise on Walk Distance in Patients with Peripheral Artery Disease (LITE). \href{https://doi.org/10.1001/jama.2021.2536}{\textit{JAMA}. 2021;325:1266-76.}
\end{enumerate}

\end{document}
